\documentclass[11pt]{article}
\usepackage[margin=1in]{geometry}
\usepackage{amsmath}
\usepackage{amssymb}
\usepackage{enumitem}
\usepackage{graphicx}
\usepackage{xcolor}
\usepackage{tcolorbox}

\title{Problem 2a: Understanding the Lagrangian \\ (For People Who Skipped Lecture)}
\author{ECO 310 - Starting from Scratch}
\date{February 23, 2026}

\begin{document}

\maketitle

\section{What Even IS a Lagrangian? (The Basics)}

\subsection{The Problem We're Trying to Solve}

Remember from Problem 1, we wanted to find the cheapest way to reach utility $u^*$. That's an \textbf{optimization problem}:

\begin{tcolorbox}[colback=blue!5!white,colframe=blue!75!black,title=The Expenditure Minimization Problem]
\textbf{Minimize:} $p_1 x_1 + p_2 x_2$ (total cost)

\textbf{Subject to:}
\begin{itemize}
\item $x_1 + x_2 \geq u^*$ (must reach utility $u^*$)
\item $x_1 \geq 0$ (can't buy negative amounts)
\item $x_2 \geq 0$ (can't buy negative amounts)
\end{itemize}
\end{tcolorbox}

\textbf{The challenge:} You can't just minimize $p_1 x_1 + p_2 x_2$ by setting everything to zero, because you have \textbf{constraints}! You need $x_1 + x_2 \geq u^*$.

\subsection{What's a Lagrangian?}

A \textbf{Lagrangian} is a mathematical tool that combines your objective (what you're minimizing/maximizing) with your constraints into ONE big equation.

\begin{center}
\textit{It's like a "master equation" that keeps track of everything at once.}
\end{center}

\textbf{Why do we need it?} Because solving constrained optimization problems directly is hard. The Lagrangian method gives us a systematic way to find the answer.

\subsection{The Magic of Lambda ($\lambda$)}

The Lagrangian uses mysterious Greek letters called \textbf{Lagrange multipliers} (written as $\lambda_1, \lambda_2, \lambda_3$, etc.).

\textbf{What are they?}
\begin{itemize}
\item They're variables (like $x$ and $y$) that we'll solve for
\item Each $\lambda$ is attached to one constraint
\item They measure "how much does the objective change if we relax this constraint slightly?"
\end{itemize}

\textbf{Intuition:} Think of each $\lambda$ as a "price" or "penalty" for violating a constraint.

\section{Breaking Down THIS Specific Lagrangian}

The problem gives you this Lagrangian:

\begin{equation}
\mathcal{L} = p_1 x_1 + p_2 x_2 + \lambda_1(-x_1) + \lambda_2(-x_2) + \lambda_3(x_1 + x_2 - u^*)
\end{equation}

Let's break it down piece by piece.

\subsection{Part 1: The Objective Function}

\begin{equation}
\boxed{p_1 x_1 + p_2 x_2}
\end{equation}

\textbf{What it is:} This is what we're trying to \textbf{minimize} - the total cost.

\textbf{Why it's here:} The Lagrangian always starts with the objective function (the thing you're optimizing).

\textbf{Example:}
\begin{itemize}
\item If coffee costs \$3 and tea costs \$5
\item And you buy 10 coffees and 0 teas
\item Cost = $3(10) + 5(0) = \$30$
\end{itemize}

\subsection{Part 2: Nonnegativity Constraint on Good 1}

\begin{equation}
\boxed{\lambda_1(-x_1)}
\end{equation}

\textbf{Original constraint:} $x_1 \geq 0$ (can't buy negative coffee)

\textbf{Why it's written as $-x_1$:}
\begin{itemize}
\item We need to rewrite $x_1 \geq 0$ as $-x_1 \leq 0$ (multiply both sides by -1)
\item Or equivalently: $0 - x_1 \leq 0$
\item The Lagrangian uses the form: constraint $\leq 0$
\end{itemize}

\textbf{What $\lambda_1$ means:}
\begin{itemize}
\item $\lambda_1$ is the "shadow price" of the nonnegativity constraint on good 1
\item If $\lambda_1 > 0$: the constraint is \textbf{binding} (you're at $x_1 = 0$)
\item If $\lambda_1 = 0$: the constraint is \textbf{not binding} (you're buying $x_1 > 0$)
\end{itemize}

\textbf{Intuition:} If you're forced to buy exactly 0 of good 1 (hitting the constraint), $\lambda_1$ tells you "how much would your cost decrease if you could buy a tiny bit?"

\subsection{Part 3: Nonnegativity Constraint on Good 2}

\begin{equation}
\boxed{\lambda_2(-x_2)}
\end{equation}

\textbf{Original constraint:} $x_2 \geq 0$ (can't buy negative tea)

\textbf{Why it's written as $-x_2$:} Same logic as above - rewrite $x_2 \geq 0$ as $-x_2 \leq 0$

\textbf{What $\lambda_2$ means:}
\begin{itemize}
\item $\lambda_2$ is the "shadow price" of the nonnegativity constraint on good 2
\item If $\lambda_2 > 0$: you're at $x_2 = 0$ (constraint is binding)
\item If $\lambda_2 = 0$: you're buying $x_2 > 0$ (constraint is not binding)
\end{itemize}

\subsection{Part 4: Utility Constraint}

\begin{equation}
\boxed{\lambda_3(x_1 + x_2 - u^*)}
\end{equation}

\textbf{Original constraint:} $x_1 + x_2 \geq u^*$ (must reach utility level $u^*$)

\textbf{Why it's written as $(x_1 + x_2 - u^*)$:}
\begin{itemize}
\item Rewrite $x_1 + x_2 \geq u^*$ as $x_1 + x_2 - u^* \geq 0$
\item For a $\geq$ constraint, we use the form: constraint $\geq 0$
\end{itemize}

\textbf{What $\lambda_3$ means:}
\begin{itemize}
\item $\lambda_3$ is the "shadow price" of utility
\item It measures: "how much would minimum cost increase if we required one more unit of utility?"
\item This is actually the \textbf{marginal cost of utility}
\item For expenditure minimization, $\lambda_3 > 0$ always (we're always hitting the utility constraint)
\end{itemize}

\textbf{Intuition:} If you need to be happier (higher $u^*$), you need to spend more. $\lambda_3$ tells you exactly how much more per unit of utility.

\section{Why This Form? The General Pattern}

\subsection{General Rules for Building Lagrangians}

For a minimization problem:

\begin{tcolorbox}[colback=yellow!10!white,colframe=orange!75!black,title=Lagrangian Recipe]
\textbf{Start with:} Objective function (what you're minimizing)

\textbf{For each constraint of type $g(x) \leq 0$:}
\begin{itemize}
\item Add: $+\lambda_i \cdot g(x)$
\item Where $\lambda_i \geq 0$
\end{itemize}

\textbf{For each constraint of type $h(x) \geq 0$:}
\begin{itemize}
\item Add: $+\mu_j \cdot h(x)$
\item Where $\mu_j \geq 0$
\end{itemize}
\end{tcolorbox}

\subsection{Applying to Our Problem}

\textbf{Objective:} Minimize $p_1 x_1 + p_2 x_2$

\textbf{Constraints:}
\begin{enumerate}
\item $x_1 \geq 0$ → rewrite as $-x_1 \leq 0$ → add $\lambda_1(-x_1)$
\item $x_2 \geq 0$ → rewrite as $-x_2 \leq 0$ → add $\lambda_2(-x_2)$
\item $x_1 + x_2 \geq u^*$ → rewrite as $(x_1 + x_2 - u^*) \geq 0$ → add $\lambda_3(x_1 + x_2 - u^*)$
\end{enumerate}

\textbf{Result:}
$$\mathcal{L} = p_1 x_1 + p_2 x_2 + \lambda_1(-x_1) + \lambda_2(-x_2) + \lambda_3(x_1 + x_2 - u^*)$$

\section{Complete Explanation (Part 2a Answer)}

\begin{tcolorbox}[colback=green!5!white,colframe=green!75!black,title=Answer to Problem 2a]
This Lagrangian represents the expenditure minimization problem for perfect substitutes.

\textbf{Each part represents:}

\begin{enumerate}
\item \textbf{$p_1 x_1 + p_2 x_2$:} The objective function - total expenditure we're minimizing

\item \textbf{$\lambda_1(-x_1)$:} Nonnegativity constraint on good 1 ($x_1 \geq 0$).
\begin{itemize}
\item $\lambda_1 > 0$ if we buy exactly 0 of good 1 (constraint binds)
\item $\lambda_1 = 0$ if we buy some good 1 (constraint doesn't bind)
\end{itemize}

\item \textbf{$\lambda_2(-x_2)$:} Nonnegativity constraint on good 2 ($x_2 \geq 0$).
\begin{itemize}
\item $\lambda_2 > 0$ if we buy exactly 0 of good 2 (constraint binds)
\item $\lambda_2 = 0$ if we buy some good 2 (constraint doesn't bind)
\end{itemize}

\item \textbf{$\lambda_3(x_1 + x_2 - u^*)$:} Utility constraint ($x_1 + x_2 \geq u^*$).
\begin{itemize}
\item Ensures we reach utility level $u^*$
\item $\lambda_3$ represents the marginal cost of utility (how much extra spending is needed for one more unit of utility)
\item This constraint always binds (we choose exactly $x_1 + x_2 = u^*$ to minimize cost)
\end{itemize}
\end{enumerate}

\textbf{Together:} The Lagrangian combines the objective with all constraints into one equation that can be optimized using calculus (taking derivatives).
\end{tcolorbox}

\section{Concrete Example: Coffee and Tea}

Let's see this in action with our coffee/tea example.

\textbf{Setup:}
\begin{itemize}
\item Coffee (good 1) costs \$3: $p_1 = 3$
\item Tea (good 2) costs \$5: $p_2 = 5$
\item Want utility 10: $u^* = 10$
\end{itemize}

\textbf{The Lagrangian:}
\begin{align*}
\mathcal{L} &= 3x_1 + 5x_2 + \lambda_1(-x_1) + \lambda_2(-x_2) + \lambda_3(x_1 + x_2 - 10) \\
&= 3x_1 + 5x_2 - \lambda_1 x_1 - \lambda_2 x_2 + \lambda_3(x_1 + x_2 - 10)
\end{align*}

\textbf{What we expect:}
\begin{itemize}
\item Coffee is cheaper, so buy only coffee: $x_1 = 10, x_2 = 0$
\item Since $x_1 = 10 > 0$, the constraint $x_1 \geq 0$ is NOT binding → $\lambda_1 = 0$
\item Since $x_2 = 0$, the constraint $x_2 \geq 0$ IS binding → $\lambda_2 > 0$
\item Since we need exactly utility 10, the utility constraint IS binding → $\lambda_3 > 0$
\end{itemize}

\textbf{Economic interpretation:}
\begin{itemize}
\item $\lambda_1 = 0$: We're not constrained by $x_1 \geq 0$ (we happily buy coffee)
\item $\lambda_2 > 0$: We'd like to buy negative tea if we could (to get a refund!), but we can't, so this constraint bites
\item $\lambda_3 > 0$: If we needed higher utility, we'd have to spend more - specifically, $\lambda_3 = 3$ (the cost per unit of utility = price of coffee)
\end{itemize}

\section{Why Do We Need the Lagrangian?}

\textbf{Can't we just solve it directly?} For perfect substitutes, yes! But for more complex utility functions (Cobb-Douglas, CES, etc.), you can't easily solve by inspection.

\textbf{What the Lagrangian gives us:}
\begin{enumerate}
\item A systematic method that works for ANY utility function
\item First-order conditions (derivatives = 0) that we can solve
\item Information about which constraints are binding (from the $\lambda$ values)
\item Economic interpretations (shadow prices)
\end{enumerate}

\section{The Big Picture}

\begin{tcolorbox}[colback=blue!5!white,colframe=blue!75!black,title=Understanding Lagrangians]
\textbf{What it is:} A tool for solving constrained optimization problems

\textbf{How it works:}
\begin{enumerate}
\item Start with objective function
\item Add each constraint, multiplied by a $\lambda$ (Lagrange multiplier)
\item Take derivatives with respect to all variables ($x_1, x_2, \lambda_1, \lambda_2, \lambda_3$)
\item Set derivatives equal to zero
\item Solve the system of equations
\end{enumerate}

\textbf{What you get:}
\begin{itemize}
\item The optimal values of $x_1, x_2$ (compensated demand)
\item The shadow prices $\lambda_i$ (marginal values of relaxing each constraint)
\item Which constraints are binding (where $\lambda_i > 0$)
\end{itemize}

\textbf{For Problem 2:} You're using this Lagrangian to derive the same answer you got in Problem 1a, but now using a formal mathematical method rather than intuition.
\end{tcolorbox}

\section{Connection to Part (b) and (c)}

\textbf{Part 2b:} Will ask you to write the optimality conditions (derivatives of $\mathcal{L}$ set equal to zero) and solve them.

\textbf{Part 2c:} Will show that when you guess wrong about which constraints bind, you get an invalid solution (like negative $\lambda$ values or violated constraints).

\textbf{Key lesson:} For perfect substitutes, exactly one nonnegativity constraint will bind (you buy only one good), and the utility constraint always binds.

\section{IMPORTANT: Two Sign Conventions (Which One Are You Using?)}

There are TWO common ways to write the Lagrangian for minimization problems. They give the same answer but look different!

\subsection{Convention 1: Positive Objective (What I Used Above)}

\begin{equation}
\mathcal{L} = p_1 x_1 + p_2 x_2 + \lambda_1(-x_1) + \lambda_2(-x_2) + \lambda_3(x_1 + x_2 - u^*)
\end{equation}

Expanding:
$$\mathcal{L} = p_1 x_1 + p_2 x_2 - \lambda_1 x_1 - \lambda_2 x_2 + \lambda_3(x_1 + x_2 - u^*)$$

\textbf{This convention:} Keeps the objective positive, adds constraints with their natural signs.

\subsection{Convention 2: NEGATIVE Objective (Probably Your Homework!)}

\begin{equation}
\mathcal{L} = -p_1 x_1 - p_2 x_2 + \lambda_1 x_1 + \lambda_2 x_2 + \lambda_3(x_1 + x_2 - u^*)
\end{equation}

Or factored:
$$\mathcal{L} = -(p_1 x_1 + p_2 x_2) + \lambda_1 x_1 + \lambda_2 x_2 + \lambda_3(x_1 + x_2 - u^*)$$

\textbf{This convention:} Negates the objective function, turning minimization into maximization.

\begin{tcolorbox}[colback=red!5!white,colframe=red!75!black,title=Why the Negative Sign?]
\textbf{Mathematical trick:} Minimizing $f(x)$ is the same as maximizing $-f(x)$!

So:
\begin{itemize}
\item \textbf{Minimize} cost $(p_1 x_1 + p_2 x_2)$
\item = \textbf{Maximize} negative cost $-(p_1 x_1 + p_2 x_2)$
\end{itemize}

Some textbooks prefer this because it makes all optimization problems look like maximization.
\end{tcolorbox}

\subsection{How to Tell Which Convention Your Problem Uses}

\textbf{Check the Lagrangian given in the problem!}

\textbf{If you see:} $\mathcal{L} = p_1 x_1 + p_2 x_2 + \ldots$ → Convention 1 (positive objective)

\textbf{If you see:} $\mathcal{L} = -p_1 x_1 - p_2 x_2 + \ldots$ → Convention 2 (negative objective)

\subsection{Breaking Down Convention 2 (Negative Objective)}

Let's explain each piece of:
$$\mathcal{L} = -(p_1 x_1 + p_2 x_2) + \lambda_1 x_1 + \lambda_2 x_2 + \lambda_3(x_1 + x_2 - u^*)$$

\subsubsection{Part 1: $-(p_1 x_1 + p_2 x_2)$}

\textbf{What it is:} Negative of total cost

\textbf{Why negative?} We're converting "minimize cost" into "maximize negative cost"

\textbf{Interpretation:} This is what we're maximizing (which is equivalent to minimizing cost)

\subsubsection{Part 2: $+\lambda_1 x_1$}

\textbf{Original constraint:} $x_1 \geq 0$ (nonnegativity on good 1)

\textbf{Why positive?} With the negative objective convention:
\begin{itemize}
\item For $\geq$ constraints, add $+\lambda \cdot (\text{constraint})$
\item So: $x_1 \geq 0$ becomes $+\lambda_1 x_1$
\end{itemize}

\textbf{What $\lambda_1$ means:} (Same as before!)
\begin{itemize}
\item $\lambda_1 > 0$ if $x_1 = 0$ (constraint binds)
\item $\lambda_1 = 0$ if $x_1 > 0$ (constraint doesn't bind)
\end{itemize}

\subsubsection{Part 3: $+\lambda_2 x_2$}

\textbf{Original constraint:} $x_2 \geq 0$ (nonnegativity on good 2)

\textbf{Why positive?} Same logic: $x_2 \geq 0$ becomes $+\lambda_2 x_2$

\textbf{What $\lambda_2$ means:}
\begin{itemize}
\item $\lambda_2 > 0$ if $x_2 = 0$ (constraint binds)
\item $\lambda_2 = 0$ if $x_2 > 0$ (constraint doesn't bind)
\end{itemize}

\subsubsection{Part 4: $+\lambda_3(x_1 + x_2 - u^*)$}

\textbf{Original constraint:} $x_1 + x_2 \geq u^*$ (utility constraint)

\textbf{Same as Convention 1!} This part doesn't change between conventions.

\textbf{What $\lambda_3$ means:} Marginal cost of utility (always positive, constraint always binds)

\subsection{Side-by-Side Comparison}

\begin{center}
\begin{tabular}{|l|c|c|}
\hline
\textbf{Part} & \textbf{Convention 1} & \textbf{Convention 2} \\
\hline
Objective & $+p_1 x_1 + p_2 x_2$ & $-p_1 x_1 - p_2 x_2$ \\
\hline
$x_1 \geq 0$ & $-\lambda_1 x_1$ & $+\lambda_1 x_1$ \\
\hline
$x_2 \geq 0$ & $-\lambda_2 x_2$ & $+\lambda_2 x_2$ \\
\hline
$x_1 + x_2 \geq u^*$ & $+\lambda_3(x_1+x_2-u^*)$ & $+\lambda_3(x_1+x_2-u^*)$ \\
\hline
\end{tabular}
\end{center}

\textbf{Key point:} The signs on the nonnegativity constraints flip, but the utility constraint stays the same!

\subsection{Complete Explanation for Convention 2}

\begin{tcolorbox}[colback=green!5!white,colframe=green!75!black,title=Answer to Problem 2a (Convention 2)]
This Lagrangian represents the expenditure minimization problem written as a maximization problem.

\textbf{Each part represents:}

\begin{enumerate}
\item \textbf{$-(p_1 x_1 + p_2 x_2)$:} Negative of expenditure - we maximize this (= minimize expenditure)

\item \textbf{$\lambda_1 x_1$:} Nonnegativity constraint on good 1 ($x_1 \geq 0$).
\begin{itemize}
\item $\lambda_1 > 0$ if we buy exactly 0 of good 1 (constraint binds)
\item $\lambda_1 = 0$ if we buy some good 1 (constraint doesn't bind)
\end{itemize}

\item \textbf{$\lambda_2 x_2$:} Nonnegativity constraint on good 2 ($x_2 \geq 0$).
\begin{itemize}
\item $\lambda_2 > 0$ if we buy exactly 0 of good 2 (constraint binds)
\item $\lambda_2 = 0$ if we buy some good 2 (constraint doesn't bind)
\end{itemize}

\item \textbf{$\lambda_3(x_1 + x_2 - u^*)$:} Utility constraint ($x_1 + x_2 \geq u^*$).
\begin{itemize}
\item Ensures we reach utility level $u^*$
\item $\lambda_3$ represents the marginal cost of utility
\item This constraint always binds ($x_1 + x_2 = u^*$)
\end{itemize}
\end{enumerate}

\textbf{Economic interpretation is identical to Convention 1!} Only the signs differ.
\end{tcolorbox}

\subsection{Which Convention Should You Use?}

\textbf{USE THE ONE YOUR HOMEWORK GIVES YOU!}

\begin{itemize}
\item Look at the Lagrangian written in the problem
\item Copy the sign convention exactly
\item The economic interpretation is the same either way
\item The optimal solution $(x_1^*, x_2^*)$ is the same either way
\end{itemize}

\begin{tcolorbox}[colback=yellow!10!white,colframe=orange!75!black,title=Bottom Line]
\textbf{If your homework shows:}
$$\mathcal{L} = -p_1 x_1 - p_2 x_2 + \lambda_1 x_1 + \lambda_2 x_2 + \lambda_3(x_1 + x_2 - u^*)$$

Then use Convention 2. Everything I explained still applies - just flip the signs on the objective and nonnegativity constraints!

\textbf{The math works out to the same answer either way.}
\end{tcolorbox}

\end{document}
